%% kapa: exam=1 round=27 year=2016 period=1 form=A subject=real_estate subject_ko=부동산학원론 seq=11 number=91
%% source: Q-Net 감정평가사 기출문제 (cst00309) artlSeq=5169402 file=extracted/1008/5169402/제27회 감정평가사 1차 문제/제27회 감정평가사 1교시문제(A형).pdf
%% conversion: kordoc 4.24.51 -> tools/gen1.py (auto); review=proofread-sonnet; flags=none
\begin{question}{91}
부동산신탁에 관한 설명으로 옳지 않은 것은?
\begin{choices}
  \item 신탁이란 위탁자가 특정한 재산권을 수탁자에게 이전하거나 기타의 처분을 하고, 수탁자로 하여금 수익자의 이익 또는 특정한 목적을 위하여 그 재산권을 관리․처분하게 하는 법률관계를 말한다.
  \item 부동산신탁의 수익자란 신탁행위에 따라 신탁이익을 받는 자를 말한다.
  \item 수익자는 위탁자가 지정한 제3자가 될 수도 있다.
  \item 신탁계약은 수익자와 위탁자 간에 체결되며 투자자는 위탁자가 발행하는 수익증권을 매입함으로써 수익자가 되어 운용성과를 얻을 수 있게 된다.
  \item 수탁자는 자산운용을 담당하는 신탁회사가 될 수 있다.
\end{choices}
\end{question}
\answer{4}
